% docs/slides/preamble.tex -- 五份 deck 共享
\documentclass[aspectratio=169,11pt]{beamer}

% ==== 中文与字体 ====
\usepackage[UTF8,fontset=none]{ctex}
\usepackage{fontspec}

\IfFontExistsTF{PingFang SC}{%
  \setCJKmainfont{PingFang SC}[BoldFont=PingFang SC,ItalicFont=PingFang SC,AutoFakeSlant=0.15]
  \setCJKsansfont{PingFang SC}[ItalicFont=PingFang SC,AutoFakeSlant=0.15]
  \setCJKmonofont{PingFang SC}
}{%
  \IfFontExistsTF{Source Han Sans SC}{%
    \setCJKmainfont{Source Han Sans SC}
    \setCJKsansfont{Source Han Sans SC}
  }{%
    \IfFontExistsTF{Noto Sans CJK SC}{%
      \setCJKmainfont{Noto Sans CJK SC}
      \setCJKsansfont{Noto Sans CJK SC}
    }{%
      \setCJKmainfont{Heiti SC}
      \setCJKsansfont{Heiti SC}
    }
  }
}

\IfFontExistsTF{Helvetica Neue}{\setmainfont{Helvetica Neue}}{%
  \IfFontExistsTF{Inter}{\setmainfont{Inter}}{}}
\IfFontExistsTF{JetBrains Mono}{\setmonofont{JetBrains Mono}[Scale=0.85]}%
  {\setmonofont{Menlo}[Scale=0.85]}

% ==== 颜色 ====
\usepackage{xcolor}
\definecolor{cMain}{HTML}{1F4E79}
\definecolor{cAccent}{HTML}{E07A1F}
\definecolor{cGray}{HTML}{595959}
\definecolor{cMute}{HTML}{F4F4F4}
\definecolor{cOK}{HTML}{2E7D32}
\definecolor{cBad}{HTML}{B23A3A}
\definecolor{cWarn}{HTML}{C8A300}

% ==== 主题 ====
\usepackage{tcolorbox}
\usetheme{metropolis}
\metroset{progressbar=frametitle,numbering=fraction,block=fill}
\setbeamercolor{frametitle}{bg=cMain,fg=white}
\setbeamercolor{progress bar}{fg=cAccent,bg=cMute}
\setbeamercolor{alerted text}{fg=cAccent}
\setbeamercolor{title}{fg=cMain}

% metropolis 默认尝试 Fira Sans，未安装时 fallback 到 Latin Modern Sans。
% 这里在主题加载后显式覆盖，强制使用 macOS 上一定存在的 Helvetica Neue，
% 跟中文 PingFang 风格一致；找不到则按顺序尝试 Inter / Arial。
\IfFontExistsTF{Helvetica Neue}{%
  \setsansfont{Helvetica Neue}%
}{%
  \IfFontExistsTF{Inter}{\setsansfont{Inter}}{%
    \IfFontExistsTF{Arial}{\setsansfont{Arial}}{}}%
}

% ==== TikZ ====
\usepackage{tikz}
\usetikzlibrary{positioning,arrows.meta,fit,shapes.geometric,calc,%
  decorations.pathmorphing,backgrounds}

\tikzset{
  node-box/.style={draw=cMain,thick,rounded corners=2pt,
    minimum height=8mm,minimum width=18mm,align=center,font=\small},
  node-state/.style={draw=cMain,thick,circle,minimum size=14mm,
    align=center,font=\small},
  node-ok/.style={draw=cOK,thick,fill=cOK!10},
  node-warn/.style={draw=cWarn,thick,fill=cWarn!15},
  node-bad/.style={draw=cBad,thick,fill=cBad!10},
  % 色盲友好的状态机三态：蓝 (正常/Closed) / 橙 (告警/Open) / 灰 (中间/HalfOpen)
  % 避免 red-green 二元区分，改用 hue + lightness 双通道。
  node-state-normal/.style={draw=cMain,thick,fill=cMain!12},
  node-state-alert/.style={draw=cAccent,thick,fill=cAccent!18},
  node-state-transition/.style={draw=cGray,thick,fill=cGray!18},
  arrow/.style={-{Stealth[length=2mm]},thick,draw=cMain},
  arrow-async/.style={-{Stealth[length=2mm]},thick,dashed,draw=cGray},
  arrow-bad/.style={-{Stealth[length=2mm]},thick,draw=cBad},
  arrow-alert/.style={-{Stealth[length=2mm]},thick,draw=cAccent},
  lifeline/.style={dashed,thick,draw=cGray},
}

% ==== 代码 ====
\usepackage{listings}
\lstdefinelanguage{Go}{
  morekeywords={break,case,chan,const,continue,default,defer,else,
    fallthrough,for,func,go,goto,if,import,interface,map,package,range,
    return,select,struct,switch,type,var,nil,true,false,iota},
  morekeywords=[2]{string,int,int64,uint,uint64,bool,byte,error,context},
  sensitive=true,
  morecomment=[l]//,morecomment=[s]{/*}{*/},
  morestring=[b]",morestring=[b]`,
}
\lstdefinelanguage{LuaRedis}{
  morekeywords={and,break,do,else,elseif,end,false,for,function,if,in,
    local,nil,not,or,repeat,return,then,true,until,while},
  morekeywords=[2]{redis,call,KEYS,ARGV,tonumber,HGET,HSET,GET,SET,
    DECR,DECRBY,INCR,INCRBY,EXPIRE,PEXPIRE,ZADD,ZCARD,ZREMRANGEBYSCORE,
    EVAL,EVALSHA},
  sensitive=true,
  morecomment=[l]{--},
  morestring=[b]",morestring=[b]',
}

\lstdefinestyle{base}{
  basicstyle=\ttfamily\scriptsize,
  numbers=left,numberstyle=\tiny\color{cGray},
  keywordstyle=\color{cMain}\bfseries,
  keywordstyle=[2]\color{cAccent},
  stringstyle=\color{cOK},
  commentstyle=\color{cGray}\itshape,
  backgroundcolor=\color{cMute},
  frame=single,framerule=0pt,
  showstringspaces=false,
  breaklines=true,
  tabsize=2,
  columns=fullflexible,
}
\lstset{style=base}

% ==== 三线表与附录页码 ====
\usepackage{booktabs}
\usepackage{appendixnumberbeamer}

% ==== 页脚来源标注 ====
\newcommand{\srcnote}[1]{{\color{cGray}\scriptsize\itshape #1}}

% 关键论点框
\newcommand{\keypoint}[1]{%
  \begin{tcolorbox}[colback=cMute,colframe=cMain,boxrule=0.5pt,arc=2pt]%
  #1\end{tcolorbox}}


\title{"翻不完的订单页" = 真实卸载}
\subtitle{7000$\times$ 提升救回多少 GMV —— gomall 列表性能的业务视角}
\author{gomall · 业务侧 deck}
\date{}

\begin{document}

\maketitle

\begin{frame}{今日大纲}
\tableofcontents
\end{frame}

% ============================================================
\section{一个被忽略 5 年的数字：p95 = 15.95 秒}
% ============================================================

\begin{frame}{先看一组真实数字（gomall 压测，order 表 6M 行）}
\begin{tabular}{@{}lrrr@{}}
\toprule
版本 & RPS & p95 延迟 & 错误率 \\
\midrule
旧版 \texttt{/orders/old/list} (OFFSET 深分页)
  & \textbf{8.3}    & \textbf{15.95s} & 0\% \\
新版 \texttt{/orders/list} (游标 + 缓存)
  & \textbf{58{,}406} & \textbf{5ms}    & 0\% \\
\midrule
\textbf{倍数差} & $\sim$7{,}000$\times$ & $\sim$3{,}200$\times$ & — \\
\bottomrule
\end{tabular}
\vspace{3mm}
\small 错误率都是 0\% —— 旧版\textbf{没崩}，只是慢得让用户以为崩了。\\
HTTP 200 + 等 16 秒 = 监控看不见的"软故障"，比 5xx 还危险。
\srcnote{stressTest/REPORT.md \S 2 订单列表表格}
\end{frame}

\begin{frame}{为什么这份 deck 关心这一个接口}
\begin{itemize}
  \item 5xx 看得见、有告警、有人盯。
  \item p95 16s 的 200 OK 在 Prometheus 默认面板上\textbf{是绿色的}。
  \item 用户不会读 HTTP code，只会卸载 app。
  \item 这一组对比是 gomall 仓库里\textbf{最戏剧化的性能数据}：吞吐 7000$\times$、延迟 3200$\times$。
  \item 它来自一行 SQL：\texttt{OFFSET 1999999 LIMIT 20}。
\end{itemize}
\keypoint{业务侧学的不是怎么改这行 SQL，是\textbf{怎么发现一行 SQL 在亏钱}。}
\srcnote{repository/db/dao/order.go:78--113 旧版 DAO 全文}
\end{frame}

% ============================================================
\section{p95 = 16s 对用户意味着什么}
% ============================================================

\begin{frame}{移动端"卡死"判定的时间窗}
\begin{tikzpicture}[node distance=0mm and 0mm,every node/.style={font=\scriptsize}]
  \draw[arrow] (0,0) -- (12,0);
  \foreach \x/\lbl in {0/0s,1.5/1s,3/2s,4.5/3s,7.5/5s,10.5/7s,12/10s+}{
    \draw[cGray] (\x,-0.08) -- (\x,0.08);
    \node[below=1mm] at (\x,-0.08) {\lbl};
  }
  \node[node-box,minimum width=22mm,fill=cOK!10,draw=cOK,above=4mm of {(2.2,0)}]
    (s1) {首屏可见};
  \node[node-box,minimum width=22mm,fill=cWarn!18,draw=cWarn,above=4mm of {(6,0)}]
    (s2) {\textbf{怀疑卡死}};
  \node[node-box,minimum width=26mm,fill=cBad!10,draw=cBad,above=4mm of {(10.5,0)}]
    (s3) {\textbf{ANR / 主动杀进程}};
\end{tikzpicture}
\vspace{4mm}
\begin{itemize}
  \item \textbf{3 秒}：行业内公认的"首屏底线"，超过开始失去注意力。
  \item \textbf{5 秒}：用户开始重复点刷新（生成无效请求 = 雪崩前奏）。
  \item \textbf{10 秒}：Android 触发 ANR 弹窗，iOS Watchdog 直接 kill 进程。
  \item \textbf{15.95s}：旧版订单列表 p95 —— 5\% 的请求每次都过线。
\end{itemize}
\end{frame}

\begin{frame}{p95 是"5\%"，不是"少数派"}
\begin{itemize}
  \item p95 = 95\% 的请求快于此值 $\Leftrightarrow$ \textbf{每 20 个用户就有 1 个体感是这个值}。
  \item 月活 100 万的小型电商：5\% = 5 万人 / 月\textbf{每次进订单页都等 16 秒}。
  \item 这 5\% 通常不是"运气差"，而是\textbf{固定群体}：
  \begin{itemize}
    \item 重度用户：订单多，翻页深，每次都踩 OFFSET 的尾部。
    \item 客户端配置弱：网络 + 设备双慢，DB 延迟被放大。
    \item 售后用户：要找 30 天前的订单 = 翻第 N 页。
  \end{itemize}
  \item 谁是这 5\%？\textbf{老客 / VIP / 投诉用户} —— 流失成本最高的一群。
\end{itemize}
\srcnote{stressTest/REPORT.md 链路压测表 第 9 行 p95=15.95s}
\end{frame}

\begin{frame}{7000$\times$ / 3200$\times$ 的视觉对比}
\begin{tikzpicture}[every node/.style={font=\scriptsize}]
  % 吞吐对比
  \node[font=\small] at (3,4) {\textbf{吞吐 (RPS)}};
  \draw[fill=cBad!40,draw=cBad] (0.5,3) rectangle (0.55,3.2);
  \node[right=1mm] at (0.55,3.1) {旧版 8.3};
  \draw[fill=cOK!40,draw=cOK]  (0.5,2.3) rectangle (6,2.5);
  \node[right=1mm] at (6,2.4) {新版 58{,}406};

  % p95 对比
  \node[font=\small] at (10,4) {\textbf{p95 延迟}};
  \draw[fill=cBad!40,draw=cBad] (7.5,3) rectangle (12.5,3.2);
  \node[right=1mm] at (12.5,3.1) {旧版 15.95s};
  \draw[fill=cOK!40,draw=cOK]  (7.5,2.3) rectangle (7.503,2.5);
  \node[right=1mm] at (7.503,2.4) {新版 5ms};

  % bottom annotation
  \node[font=\footnotesize,align=center] at (6.5,1.3) {
    左图：bar 长按比例画，旧版仅 0.014\% 宽度（$8.3/58{,}406$）。\\
    右图：旧版条占满，新版条仅 0.03\% 宽度（$5\text{ms}/15{,}950\text{ms}$）。
  };
\end{tikzpicture}
\vspace{2mm}
\small 7000$\times$ 不是"快了一点"，是\textbf{把"几乎所有请求都失败"变成"几乎所有请求都成功"}。
\srcnote{stressTest/REPORT.md \S 2 订单列表代差碾压表}
\end{frame}

% ============================================================
\section{8.3 RPS 是什么概念}
% ============================================================

\begin{frame}{8.3 RPS = 一台机器只够 8 个并发}
\begin{tikzpicture}[node distance=4mm and 6mm]
  \node[node-box,node-ok,minimum width=18mm] (u1) {用户 1};
  \node[node-box,node-ok,minimum width=18mm,right=of u1] (u2) {用户 2};
  \node[node-box,node-ok,minimum width=18mm,right=of u2] (u3) {用户 3};
  \node[node-box,node-ok,minimum width=18mm,right=of u3] (u4) {用户 4};
  \node[node-box,node-ok,minimum width=18mm,right=of u4] (u5) {用户 5};
  \node[node-box,node-ok,minimum width=18mm,right=of u5] (u6) {\dots 8};
  \node[node-box,node-bad,minimum width=18mm,right=of u6] (u7) {\textbf{用户 9}};
  \draw[arrow-bad,bend left=20] (u7.north) to node[above,font=\tiny]{超时 / 白屏} (u7.north east);
\end{tikzpicture}
\vspace{2mm}
\begin{itemize}
  \item 8.3 RPS 意味着\textbf{每秒只能服务 8 个新请求}。
  \item 运营在企业微信群发"双 11 我的订单"链接 $\to$ 群里 200 人点 $\to$ 99\% 看到的是\textbf{加载圈圈}。
  \item 推送通知尤其危险：\textbf{瞬时并发是平均值的 10--100 倍}。
  \item 客服后台批量查"为什么这个用户说没收到货" $\to$ 一查一个超时。
\end{itemize}
\srcnote{stressTest/REPORT.md 链路压测表 第 9 行 RPS=8.3 / VUs=100}
\end{frame}

\begin{frame}{相同的反模式还藏在 \texttt{product/list}}
\begin{tabular}{@{}lrr@{}}
\toprule
接口 & RPS & p95 \\
\midrule
\texttt{/orders/old/list} （\texttt{OFFSET 1999999}） & 8.3  & 15.95s \\
\texttt{/product/list}    （\texttt{COUNT(*)} 全表）  & 24.5 & 2.50s \\
\bottomrule
\end{tabular}
\vspace{3mm}
\begin{itemize}
  \item \texttt{product/list} 没有 OFFSET 病灶，但每次都 \texttt{SELECT COUNT(*)} 956K 行的 product 表。
  \item COUNT 走全表 = 单次 2.5s，瓶颈在\textbf{要不要告诉用户"共 X 页"}。
  \item 这是\textbf{产品决策}：要 "共 X 页" 还是要 24.5 $\to$ 5{,}000 RPS。
  \item gomall 现状：还没改 —— 这是仓库里\textbf{第二大的性能债}。
\end{itemize}
\srcnote{repository/db/dao/product.go:59--64 \texttt{CountProductByCondition}}
\end{frame}

% ============================================================
\section{为什么旧版是 OFFSET 1999999 LIMIT 20}
% ============================================================

\begin{frame}[fragile]{旧版 SQL 的真实样子}
\begin{lstlisting}[language=Go,firstnumber=85,
  caption={\srcnote{repository/db/dao/order.go:85--105}}]
query := dao.DB.Table("order").Where(
    "`order`.user_id = ? and `order`.type = ?", uId, req.Type)

query.Count(&count)              // 每页都 COUNT(*) 全表
offset := (req.PageNum - 1) * req.PageSize
err = query.
    Joins("left join product as p on p.id=order.product_id").
    Joins("left join address as a on a.id=order.address_id").
    Order("order.created_at desc").   // 非主键排序 = filesort
    Offset(offset).                   // 翻第 N 页 = 扫前 N*20 行
    Limit(req.PageSize).
    Find(&r.List).Error
\end{lstlisting}
\small 三个性能病灶叠在 20 行里：COUNT 全表 + filesort + OFFSET。
\end{frame}

\begin{frame}{OFFSET 1999999 LIMIT 20 是怎么扫表的}
\begin{tikzpicture}[node distance=2mm and 4mm,every node/.style={font=\scriptsize}]
  % 一行行扫
  \foreach \i/\x in {1/0, 2/0.6, 3/1.2, 4/1.8, 5/2.4, 50000/3.2, 1999999/6.0}{
    \node[draw=cBad,fill=cBad!15,minimum width=8mm,minimum height=4mm] at (\x,0) {扫};
  }
  \node[draw=cOK,fill=cOK!30,minimum width=10mm,minimum height=4mm] at (7.0,0) {返回};
  \draw[arrow-bad] (-0.5,0) -- node[above]{依次扫 1{,}999{,}999 行（必须真读）} (5.7,0);
  \draw[arrow] (6.5,0) -- node[above]{只要这 20 行} (8,0);

  \node[font=\scriptsize,align=left] at (4.5,-1.2) {
    数据库\textbf{不能}"跳到第 200 万行" —— 它必须从 id=1 开始一行一行扫，\\
    扔掉前 1{,}999{,}999 行，才知道第 1{,}999{,}999 行长什么样。\\
    JOIN product + address + filesort 全在这 200 万行上发生。
  };
\end{tikzpicture}
\srcnote{repository/db/dao/order.go:91--104 \texttt{Offset} 行为}
\end{frame}

\begin{frame}{为什么会写成这样：因为前端要"第 N 页"}
\begin{itemize}
  \item 早期产品原型来自\textbf{后台管理界面}：表格 + "上一页 / 下一页 / 跳转到 X 页"。
  \item 这个交互假设的是 \textbf{运营 / 客服}，不是\textbf{消费者}。
  \item 消费者根本不会跳到第 100 页查订单 —— 他要么找最近的，要么按订单号搜。
  \item \textbf{设计错位}：把后台的分页直接搬到 C 端。
  \item 修复方案不只是"换一行 SQL"，是\textbf{产品交互的重构}：
  \begin{itemize}
    \item C 端订单列表 $\to$ 无限滚动 / 下拉加载。
    \item 历史订单查询 $\to$ 按订单号 / 时间范围搜索。
    \item 客服后台保留分页，但限制最大页码。
  \end{itemize}
\end{itemize}
\srcnote{routes/router.go:78 新接口 \texttt{/orders/list}；:79 旧接口 \texttt{/orders/old/list}}
\end{frame}

\begin{frame}{新版 SQL 的样子：游标 + LastId}
\begin{tikzpicture}[node distance=4mm and 6mm]
  \node[node-box,node-ok] (req) {第 1 页\\\scriptsize last\_id = 0};
  \node[node-box,right=of req] (sql1) {\texttt{ORDER BY id DESC}\\\texttt{LIMIT 20}};
  \node[node-box,right=of sql1,node-warn] (cache) {Redis 5min};
  \node[node-box,below=of req] (req2) {第 N 页\\\scriptsize last\_id = 上一页末尾};
  \node[node-box,right=of req2] (sql2) {\texttt{WHERE id < last\_id}\\\texttt{LIMIT 20}};
  \node[node-box,right=of sql2,node-ok] (idx) {PK 索引\\O(1) 命中};
  \draw[arrow] (req)  -- (sql1);
  \draw[arrow] (sql1) -- (cache);
  \draw[arrow] (req2) -- (sql2);
  \draw[arrow] (sql2) -- (idx);
\end{tikzpicture}
\vspace{3mm}
\begin{itemize}
  \item 不存在"跳到第 200 万行"这件事 —— 客户端只能"从这里继续往下"。
  \item 数据库每次只扫 20 行：从 last\_id 沿 PK 索引倒序拿。
  \item 第一页结果再加 5min Redis 缓存 $\to$ 多数请求不进 DB。
\end{itemize}
\srcnote{repository/db/dao/order.go:50--75 新版 DAO 全文}
\end{frame}

% ============================================================
\section{7000$\times$ 救回了什么}
% ============================================================

\begin{frame}{救回的不是 RPS，是三条业务损失}
\begin{tabular}{@{}lll@{}}
\toprule
角度 & 旧版亏的是 & 新版救回的是 \\
\midrule
用户体验 & 16s 白屏 $\to$ ANR     & 3 秒内可见 \\
转化漏斗 & 未付款用户看不到订单 & 翻回订单 $\to$ 继续支付 \\
客服压力 & "为什么订单不显示"   & 投诉降到接近零 \\
\bottomrule
\end{tabular}
\vspace{3mm}
\begin{itemize}
  \item 第一条 = \textbf{获客成本浪费}：广告砸来的用户在订单页流失。
  \item 第二条 = \textbf{GMV 直接损失}：下了单不付 $\to$ 占库存 $\to$ 30 分钟后 Saga 释放 $\to$ 真正想买的人那时已经走了。
  \item 第三条 = \textbf{客服成本}：一通客服电话 $\approx$ 8 元，月 5 万投诉 $\approx$ 40 万 / 月。
\end{itemize}
\srcnote{service/order.go:104--122 新版；:124--142 旧版的 \texttt{OrderListOld}}
\end{frame}

\begin{frame}{"看不到订单 $\to$ 不能付款"的漏斗损失}
\begin{tikzpicture}[node distance=6mm and 10mm]
  \node[node-box,node-ok] (s1) {下单成功\\\scriptsize 100\%};
  \node[node-box,node-warn,right=of s1] (s2) {跳转到订单列表\\\scriptsize 95\% 进来};
  \node[node-box,node-bad,right=of s2] (s3) {等 16s 白屏\\\scriptsize -30\% 流失};
  \node[node-box,node-warn,below=of s3] (s4) {点支付\\\scriptsize 剩 65\%};
  \node[node-box,node-ok,left=of s4] (s5) {支付成功\\\scriptsize $\approx$ 58\%};
  \draw[arrow]      (s1) -- (s2);
  \draw[arrow-bad]  (s2) -- (s3);
  \draw[arrow-bad]  (s3) -- (s4);
  \draw[arrow]      (s4) -- (s5);
\end{tikzpicture}
\vspace{2mm}
\begin{itemize}
  \item 行业 baseline：下单 $\to$ 付款转化在 90\%+。
  \item 16s 白屏挤掉一条边 $\to$ 整条漏斗少 30 个百分点。
  \item 假设客单价 200，月活 100 万，新版救回的 30\% $\approx$ \textbf{600 万 GMV / 月}。
\end{itemize}
\small 数字是简化模型，但量级方向是对的：性能问题首先是\textbf{钱的问题}。
\end{frame}

\begin{frame}{客服侧：投诉从"加载失败"消失}
\begin{itemize}
  \item 旧版客诉 top 3 关键词（业界典型）：
  \begin{itemize}
    \item "订单页一直转圈"
    \item "为什么我的订单不显示"
    \item "刚买的东西在哪里"
  \end{itemize}
  \item 这三类\textbf{本质都是同一个 bug}：列表太慢用户以为没下单 $\to$ 重新下单 $\to$ 双单 $\to$ 客诉。
  \item 双单还会触发幂等\textbf{无关的}场景：用户主动点两次 = 业务层重复，幂等中间件\textbf{挡不住}（不同 idempotency key）。
  \item 新版上线后，这一类\textbf{从客诉报表中整类消失}。
  \item 客服培训成本、退款工单成本、商家结算成本都跟着降。
\end{itemize}
\srcnote{deck 01 idempotency 解决的是请求级重放，不解决"用户主动点两次"}
\end{frame}

% ============================================================
\section{谁来发现这种 bug}
% ============================================================

\begin{frame}{开发的本地环境永远看不到这个 bug}
\begin{tikzpicture}[node distance=6mm and 10mm]
  \node[node-box,node-ok] (dev) {本地开发\\\scriptsize 20 行测试数据};
  \node[node-box,node-warn,right=of dev] (uat) {UAT\\\scriptsize 1 万行};
  \node[node-box,node-bad,right=of uat] (prod) {生产\\\scriptsize 6M 行 order};
  \draw[arrow] (dev) -- node[above,font=\tiny]{快} (uat);
  \draw[arrow-bad] (uat) -- node[above,font=\tiny]{慢 100$\times$} (prod);
\end{tikzpicture}
\vspace{3mm}
\begin{itemize}
  \item 20 行测试数据：OFFSET 1{,}999{,}999 根本跑不到 —— 全表才 20 行。
  \item 1 万行 UAT：能跑到但延迟在 100ms 级，测试同学不会标红。
  \item 6M 行生产：OFFSET 走索引也能扫满磁盘 IO，p95 飙到 16s。
  \item \textbf{结论：本地单测 / UAT 都看不见这个 bug}。
\end{itemize}
\keypoint{性能 bug 不是"在开发机能复现"的 bug，必须有\textbf{大表压测}。}
\srcnote{stressTest/REPORT.md 文件头"数据规模 order 表 6M 行"}
\end{frame}

\begin{frame}{真实发现链路：客服 $\to$ RUM $\to$ 慢日志}
\begin{tikzpicture}[node distance=8mm and 14mm]
  \node[node-box] (cs) {客服};
  \node[node-box,right=of cs,node-warn] (rum) {RUM\\\scriptsize 前端埋点};
  \node[node-box,right=of rum,node-warn] (apm) {APM\\\scriptsize p95 告警};
  \node[node-box,right=of apm,node-bad] (sql) {慢查询日志\\\scriptsize MySQL};
  \node[node-box,below=of apm] (dev) {开发};
  \draw[arrow] (cs)  -- (rum);
  \draw[arrow] (rum) -- (apm);
  \draw[arrow] (apm) -- (sql);
  \draw[arrow] (rum) -- (dev);
  \draw[arrow] (sql) -- (dev);
\end{tikzpicture}
\vspace{2mm}
\begin{itemize}
  \item \textbf{客服}：第一反应是"工具坏了"，能拿到\textbf{用户 ID + 时间窗}。
  \item \textbf{RUM}（Real User Monitoring）：埋点 \texttt{page\_load\_time}，看真实用户分布。
  \item \textbf{APM}：从 trace 里看 p95，能定位到\textbf{接口}。
  \item \textbf{慢查询日志}：定位到\textbf{具体 SQL}，再交给开发。
  \item 单靠 APM 平均值看不见 p95 异常 —— 必须\textbf{看分位数 + RUM 真实分布}。
\end{itemize}
\end{frame}

\begin{frame}{SLO 告警的设计：监控分位数而非平均值}
\begin{tabular}{@{}lll@{}}
\toprule
指标 & 旧版实测 & 内部 SLO 应当 \\
\midrule
avg 延迟       & 10.5s   & p99 < 500ms \\
p95 延迟       & 15.95s  & p95 < 200ms \\
错误率（HTTP）  & 0\%     & < 0.1\% \\
\textbf{告警是否会响} & \textbf{不会} & \textbf{应该响} \\
\bottomrule
\end{tabular}
\vspace{3mm}
\begin{itemize}
  \item 旧版"错误率 0\%"是 gomall 监控的盲区：HTTP 200 就算成功。
  \item 应当加 \textbf{延迟分位数 SLO}：p95 > 1s 持续 5min 报警。
  \item RUM 侧再加\textbf{用户感知 SLO}：首屏 > 3s 占比超过 10\% 报警。
  \item 两者交叉验证，才能拦住"HTTP 200 + 等 16 秒"这种软故障。
\end{itemize}
\srcnote{stressTest/REPORT.md \S 2 旧版 p50=10.5s / p95=15.95s / max=16.29s}
\end{frame}

% ============================================================
\section{四类角色看慢查询}
% ============================================================

\begin{frame}{产品 / 运营 / 客服 / SRE 关心什么}
\begin{tabular}{@{}llll@{}}
\toprule
角色 & 关心 & 看哪个指标 & 旧版的回答 \\
\midrule
产品  & 转化率 / 跳出率   & 漏斗 / 首屏 LCP        & 砍 30 个点 \\
运营  & 大促 / 推送承受  & 峰值 RPS / 误差率      & 8.3 RPS \\
客服  & 工单类型分布     & 关键词 "加载" / "卡"    & 月 5 万投诉 \\
SRE  & 可用率 / 容量    & p95 / p99 / error rate & 200 OK 但 16s \\
\bottomrule
\end{tabular}
\vspace{2mm}
\begin{itemize}
  \item 四个角色看的是\textbf{同一个 bug 的不同切面}。
  \item 单看任何一个都不会触发 P0 故障 —— 必须\textbf{把四个切面拉到一张表}。
  \item gomall 现状：四个面都看不见，因为没有 RUM / 没有客诉关键词分类 / 监控没看 p95。
\end{itemize}
\srcnote{routes/router.go:78--79 同一份业务两个 endpoint，方便 AB 对比}
\end{frame}

\begin{frame}{运营视角：推送链接秒杀掉 8 RPS 接口}
\begin{itemize}
  \item 运营做"双 11 我的订单一键查询"活动 $\to$ APP 推送 200 万 DAU。
  \item 推送打开率 5\% $\to$ 10 万人在\textbf{1 分钟内}点开。
  \item 旧版 8.3 RPS = 60s $\times$ 8.3 = \textbf{500 个请求能服务，剩下 99{,}500 个超时}。
  \item 运营复盘："为什么转化这么差" $\to$ 看不出问题，因为\textbf{服务器没崩}。
  \item 这是为什么"一个接口慢"在运营看来是\textbf{活动效果差}，看不到接口层。
\end{itemize}
\keypoint{运营做活动前必须问后端："这接口能扛多少 RPS？" —— gomall 现在的答案应该是 58K，不是 8.3。}
\srcnote{stressTest/REPORT.md \S 2 RPS 7000$\times$ 表格}
\end{frame}

% ============================================================
\section{反思：为什么开发会写出这种代码}
% ============================================================

\begin{frame}{常见误判：5 个让开发踩这个坑的理由}
\begin{enumerate}
  \item "GORM 的 Offset/Limit 就是标准分页姿势" —— 标准但\textbf{不可扩展}。
  \item "PageNum + PageSize 是 RESTful 习惯" —— REST 没规定一定 OFFSET。
  \item "前端要 total 显示 X 页" —— 前端没想过\textbf{近似 total} 是可接受的。
  \item "本地测试都好的" —— 因为本地没有 6M 行。
  \item "线上没人投诉" —— 因为 HTTP 200 没有触发监控。
\end{enumerate}
\vspace{3mm}
\small 每条都是\textbf{独立的合理判断}，叠在一起就是 16 秒的 p95。\\
所以代码审查很难抓住 —— reviewer 看的是\textbf{这一行}，不是\textbf{这一行在 6M 数据下的行为}。
\srcnote{repository/db/dao/order.go:80--105 旧版的注释直接写出三处性能病灶}
\end{frame}

\begin{frame}{制度建设：怎么不再次踩}
\begin{itemize}
  \item \textbf{大表压测进 CI}：每个列表接口必须跑 \texttt{stressTest/} 下脚本，p95 > 1s 红线。
  \item \textbf{评审 checklist 加一条}："这个接口在 100 万行下 p95 多少？" 没数 $\to$ 拒 merge。
  \item \textbf{监控强制 p95 告警}：任何 P1 接口的 p95 > 500ms 持续 5min 报警。
  \item \textbf{产品 review 加一条}：列表设计先选\textbf{无限滚动 / 游标}，要 PageNum 的写 RFC 解释。
  \item \textbf{数据库治理}：用 \texttt{pt-query-digest} 周报，把 OFFSET > 1000 的 SQL 揪出来。
\end{itemize}
\srcnote{stressTest/order\_list\_deep\_pagination.js 和 order\_list\_optimized.js 已经成对}
\end{frame}

% ============================================================
\section{回到业务：性能即业务}
% ============================================================

\begin{frame}{把这一份 deck 浓缩成一句话}
\begin{itemize}
  \item 旧版没崩，所以监控没响。
  \item 监控没响，所以 SRE 不知道。
  \item SRE 不知道，所以产品不知道。
  \item 产品不知道，所以"为什么转化下降" 找不到答案。
  \item 答案是：\textbf{一行 SQL 让 5\% 的用户每次等 16 秒}。
\end{itemize}
\keypoint{性能问题不是技术问题，是\textbf{业务问题}：它直接砸 GMV、扩客服编制、消耗市场预算。}
\srcnote{repository/db/dao/order.go:78--113 旧版 / :35--76 新版的一行 SQL 改写}
\end{frame}

% ============================================================
\appendix
% ============================================================

\begin{frame}{业务 Q\&A（上）}
\textbf{Q1}：为什么不用普通的 \texttt{PageNum + PageSize}？\\
\hspace{4mm}\small A：因为电商订单列表用户根本不翻第 100 页，但 OFFSET 让每翻一页都扫前 N$\times$20 行。\textbf{交互上} 99\% 用户用不到的"跳页"功能，\textbf{代价是}所有用户每次都付。

\vspace{2mm}
\textbf{Q2}：怎么发现这种慢的接口？\\
\hspace{4mm}\small A：客服关键词 + RUM 首屏 LCP + APM p95 + 慢查询日志。四条线交叉。单靠错误率看不见 —— 旧版错误率 0\%。

\vspace{2mm}
\textbf{Q3}：能不能折中：保留 PageNum 但只允许翻前 10 页？\\
\hspace{4mm}\small A：可以。客服后台保留分页（最大 100 页）。C 端走游标。产品需求 = 用户最近 N 单 + 搜索旧单，覆盖 99\%。
\end{frame}

\begin{frame}{业务 Q\&A（下）}
\textbf{Q4}：\texttt{product/list} 的 2.5s 怎么办？是不是和订单一样换游标？\\
\hspace{4mm}\small A：不是。\texttt{product/list} 的瓶颈是 \texttt{COUNT(*) }全表，不是 OFFSET。三个方向：(a) 用近似 count（\texttt{SHOW TABLE STATUS}）；(b) 给 \texttt{category\_id} 建索引让 COUNT 走索引；(c) 产品上不再显示"共 X 页"。是\textbf{产品 + 开发}一起做的决定。

\vspace{2mm}
\textbf{Q5}：救回 30\% 转化这个数字怎么算的？\\
\hspace{4mm}\small A：模型简化。但量级方向看：行业基线下单$\to$支付转化 90\%+，16s 白屏挤掉的部分一定是两位数百分比。具体数字以\textbf{上线前后 AB 测}的实测为准 —— 这也是 \texttt{/orders/list} 和 \texttt{/orders/old/list} 同时保留的原因。

\vspace{2mm}
\textbf{Q6}：为什么开发不在本地就测出来？\\
\hspace{4mm}\small A：本地表里几十行 $\to$ OFFSET 1{,}999{,}999 等于 LIMIT 0。必须\textbf{大表压测进 CI}。gomall 已经有 \texttt{stressTest/} 一套脚本，下一步该接入 GitHub Actions。
\end{frame}

\begin{frame}{代码位置一览}
\begin{description}
  \item[路由：新 / 旧两个入口]\srcnote{routes/router.go:78--79}
  \item[新版游标分页 DAO]\srcnote{repository/db/dao/order.go:35--76}
  \item[旧版 OFFSET 深分页 DAO]\srcnote{repository/db/dao/order.go:78--113}
  \item[订单 service 新版]\srcnote{service/order.go:104--122}
  \item[订单 service 旧版]\srcnote{service/order.go:124--142}
  \item[订单列表请求结构]\srcnote{types/order.go:30--46（LastId 字段）}
  \item[product COUNT(*) 全表]\srcnote{repository/db/dao/product.go:59--64}
  \item[压测脚本对]\srcnote{stressTest/order\_list\_deep\_pagination.js + order\_list\_optimized.js}
  \item[压测报告]\srcnote{stressTest/REPORT.md \S 2 / \S 3}
\end{description}
\vspace{2mm}
\keypoint{配套技术 deck：03-cache-consistency.tex \S 游标分页 + 首页缓存。}
\end{frame}

\end{document}
